\chapter{From Angles and Distance to X, Y, Z}

\section{The idea}

\begin{kid}
To tell a friend where a treasure is hidden, you could say: ``Face the big tree. Turn a little to the left. Look a little up. Now walk 25 steps that way.'' That is a \textbf{turn}, a \textbf{tilt} and a \textbf{distance}. A LiDAR records the same three things for every dot. Then it does some triangle math to change them into ``how far forward, how far sideways, how far up'' --- which we call X, Y and Z.
\end{kid}

\fig[0.6]{polar3d}{A range of 25 m with horizontal angle 40° and vertical angle 15° becomes a point at (18.50, 15.52, 6.47) m.}

\fig[0.95]{polar_steps}{The two moves: tilt (side view), then turn (top view).}

\section{Mechanical versus solid-state scanners}

\fig[1.0]{patterns}{How three kinds of scanner sample the world (angles in degrees).}

\section{How dense will my points be?}

\begin{kid}
Imagine painting a road with a paint-roller that only touches the road once every big step. If you walk faster, the stripes get further apart. If you roll faster, they get closer.
\end{kid}

\fig[0.85]{spacing}{Each sweep of the scanner lays a line of points across the road.}

\begin{tip}
Mobile mapping heads for survey work are often mounted \emph{tilted} so the spinning plane slices the road at an angle; then several channels hit the same strip of road at different times and along-track gaps close up. Point-density specifications should be checked on the delivered data (points per m$^2$ on a test patch), not assumed from the brochure.
\end{tip}
